\section{Identificação do Projeto}

\subsection{Projeto e contexto académico}

O presente documento descreve a organização do repositório do projeto \textbf{NatureProtector}, desenvolvido no contexto da unidade curricular de Projeto e Seminário da Licenciatura em Engenharia Informática e de Computadores do Instituto Superior de Engenharia de Lisboa.

O NatureProtector é um sistema de apoio à prevenção e monitorização operacional de risco de incêndio rural. O repositório integra o código fonte, os testes, a infraestrutura local, a interface web, os scripts de execução, os dados de suporte, a documentação técnica, a evidência recolhida durante a validação e os materiais de entrega associados ao projeto.

\begin{longtable}{L{0.30\textwidth} L{0.60\textwidth}}
\toprule
\textbf{Campo} & \textbf{Descrição} \\
\midrule
\endfirsthead
\toprule
\textbf{Campo} & \textbf{Descrição} \\
\midrule
\endhead
Projeto & NatureProtector \\
Instituição & Instituto Superior de Engenharia de Lisboa \\
Curso & Licenciatura em Engenharia Informática e de Computadores \\
Unidade curricular & Projeto e Seminário \\
Autores & Miguel Alves; Gabriel Mano \\
Orientadores & Artur Ferreira; Nuno Leite \\
Repositório & \href{https://github.com/MAlves4018/NatureProtector.git}{GitHub público do projeto} \\
\bottomrule
\end{longtable}

\subsection{Finalidade do documento}

A finalidade deste documento é apresentar a organização do repositório de forma objetiva, permitindo a um avaliador, orientador ou novo elemento da equipa localizar rapidamente os principais artefactos do projeto.

O documento identifica:

\begin{itemize}
    \item a localização pública do repositório;
    \item a estrutura de alto nível das pastas;
    \item os projetos de código fonte e respetivas responsabilidades;
    \item os projetos de teste;
    \item a interface web;
    \item os dados, manifestos, scripts, infraestrutura e documentação;
    \item os comandos principais de validação;
    \item os artefactos gerados que não devem ser confundidos com código fonte.
\end{itemize}

\subsection{Âmbito e limites}

Este documento não substitui o relatório técnico do projeto, a documentação de arquitetura nem o guia completo de setup. A sua função é descrever a organização do repositório e indicar onde se encontram os artefactos relevantes.

Assim, ficam fora do âmbito deste documento:

\begin{itemize}
    \item a descrição científica completa da fórmula de risco;
    \item a fundamentação detalhada dos índices FWI e KBDI;
    \item a discussão metodológica completa da V1;
    \item a validação científica final do modelo;
    \item o manual exaustivo de utilização da aplicação.
\end{itemize}

A referência principal para preparar o ambiente local é o ficheiro \path{docs/local-baseline-setup.md}. Este documento apenas indica essa referência e resume os comandos de validação mais relevantes.
